\section{Specializace vrstev: čtení labelů a kernelové porovnávání}
\label{sec:specializace-vrstev}

V předchozích dvou sekcích jsme se pohybovali výhradně mezi šestivrstvými modely – nejprve fixní model z~oddílu~\ref{sec:vnitrni-struktura}, poté náhodný model z~oddílu~\ref{sec:pfn-random-hp}.
Ukázali jsme, že PFN neprovádí prostý kernel smoothing a že náhodný model dokáže identifikovat korelační délku z~kontextu.
Zůstává však otevřená otázka \emph{jak} tento výpočet probíhá uvnitř transformeru – co dělá která vrstva a proč jich model vlastně potřebuje víc než jednu.
V této sekci proto poprvé systematicky měníme \textbf{hloubku} sítě a ptáme se, zda se mezi vrstvami vytváří dělba práce.
\par

Analýzu provádíme na rodině pěti modelů, které se liší \textit{pouze} počtem vrstev: $1$, $2$, $4$, $6$ a $8$ vrstev.
Všechny sdílejí identickou architekturu s modely z~předchozích sekcí ($8$ attention hlav, dimenze embeddingu $512$, dimenze skryté vrstvy dopředné sítě $1024$, \lstinline|features_per_group=1| a výstupní \texttt{Full\-Support\-Bar\-Distribution} s~$1000$ biny) a všechny byly trénovány na téže distribuci hyperparametrů jako náhodný model:
$$\ell \sim \mathcal{U}(0{,}05,\, 1{,}0), \qquad \text{outputscale} \sim \text{LogNormal}(0,\, 0{,}5), \qquad \sigma^2 \sim 10^{\mathcal{U}(-3,\,-1)}.$$
Počet parametrů roste s~hloubkou od $3{,}65$ milionu u~jednovrstvého modelu po $14{,}14$ milionu u~šestivrstvého.
Jednovrstvý model konvergoval již za $100$ epoch, hlubší modely byly trénovány $300$ epoch.
Tato kontrolovaná řada nám umožňuje sledovat, jak se vnitřní strategie transformeru proměňuje s~přidáváním vrstev.
\par

Motivace pro tuto analýzu vychází přímo z~Naglerovy teoretické práce.
Připomeňme, že v~PFN je každý kontextový bod reprezentován tokenem $V_j = (Y_j,\, X_j)^\top$, který v~sobě slučuje jak vstupní souřadnici $X_j$, tak její pozorovanou hodnotu $Y_j$.
Attention skóre jednovrstvého modelu je proto \textit{lineární kombinací obojího} – zhruba tvaru $\alpha X_j + \beta Y_j$ – na rozdíl od GP váhy $k(x^*, X_j)$, která závisí \textit{výhradně} na vstupní souřadnici.
Nagler odtud odvozuje, že jediná vrstva nemůže reprodukovat GP posterior přesně: tzv.\ tilted-measure limit se nikdy nezlokalizuje kolem $x^*$, a proto zůstává v~predikci nenulový bias.
To přirozeně vede k~hypotéze o~dělbě práce: pokud jedna vrstva nestačí, možná se hlubší modely naučí \textit{rozdělit} úkol – rané vrstvy by mohly číst hodnoty $Y_j$ (label reading) a prostřední vrstvy porovnávat vzdálenosti ve~vstupním prostoru $X$ (kernel-like chování).
Právě tuto hypotézu v~následujících experimentech testujeme.
\par

\subsection{Korelační analýza attention napříč vrstvami}
\label{subsec:ch2-korelace}

Abychom mohli změřit, „na co se která vrstva dívá“, spočítáme pro každou vrstvu a každou z~jejích osmi hlav dvě korelace attention vah:
korelaci s~\textbf{RBF/NW kernelem} $k(x^*, X_j)$, která měří, nakolik se hlava chová jako kernelové porovnání vstupních vzdáleností, a korelaci s~\textbf{labely} $Y_j$, která měří, nakolik hlava místo toho čte pozorované hodnoty.
Experiment provádíme na $200$ instancích se sdílenými daty pro všechny modely, s~kontextem $n_{\text{ctx}} = 40$, $n_{\text{test}} = 20$ a referenčním kernelem s~$\ell_{\text{NW}} = 0{,}5$.
\par

\textbf{Proč Spearman a ne Pearson:}

Attention váhy jsou silně zešikmené na simplexu – naprostá většina hmoty leží na několika málo bodech a zbytek je téměř nulový.
Pearsonova korelace, citlivá na absolutní hodnoty, proto monotónní vztah mezi attention a kernelem systematicky podhodnocuje.
Přechodem na \textit{Spearmanovu} (pořadovou) korelaci se kernelové korelace ve~středních vrstvách zvedly z~přibližně $0{,}35$ na $\sim 0{,}8$ – trend zůstává týž, ale je mnohem výraznější.
V~celé této sekci proto uvádíme Spearmanovy korelace.
\par

Přehled kernelových korelací zprůměrovaných přes hlavy shrnuje tabulka~\ref{tab:ch2_corr_overview} a graficky jej spolu s~labelovými korelacemi zachycuje obrázek~\ref{fig:ch2_overview}.

\begin{table}[htbp]
\centering
\begin{tabular}{c|ccccc}
\hline
\textbf{Vrstva} & \textbf{1 vrstva} & \textbf{2 vrstvy} & \textbf{4 vrstvy} & \textbf{6 vrstev} & \textbf{8 vrstev} \\
\hline
$0$ & $\mathbf{0{,}480}$ & $0{,}168$ & $0{,}074$ & $0{,}004$ & $-0{,}099$ \\
$1$ & — & $\mathbf{0{,}819}$ & $\mathbf{0{,}887}$ & $0{,}610$ & $0{,}563$ \\
$2$ & — & — & $0{,}828$ & $\mathbf{0{,}826}$ & $\mathbf{0{,}873}$ \\
$3$ & — & — & $0{,}270$ & $0{,}817$ & $0{,}765$ \\
$4$ & — & — & — & $0{,}740$ & $0{,}773$ \\
$5$ & — & — & — & $0{,}416$ & $0{,}544$ \\
$6$ & — & — & — & — & $0{,}530$ \\
$7$ & — & — & — & — & $0{,}296$ \\
\hline
\end{tabular}
\caption{Spearmanova korelace attention vah s~RBF/NW kernelem, zprůměrovaná přes $8$ hlav dané vrstvy a přes $200$ instancí. Pomlčka značí, že daná vrstva u~mělčího modelu neexistuje. Tučně je vyznačeno maximum ve~sloupci. Ve~vrstvě $0$ je kernelová korelace vždy nízká (u~$8$ vrstev dokonce mírně záporná), zatímco ve~středních vrstvách skáče přes $0{,}8$.}
\label{tab:ch2_corr_overview}
\end{table}

\begin{figure}[p]
    \centering
    \includegraphics[width=0.92\linewidth]{figures/ch2_label_kernel/fig_ch2_q1_overview_all_models.png}
    \caption{Srovnání všech pěti modelů. \textbf{Nahoře:} Spearmanova korelace attention vah s~RBF kernelem jako funkce indexu vrstvy. \textbf{Dole:} absolutní hodnota korelace attention vah s~labely $Y_j$. Barevné pásy značí $\pm 1$ směrodatnou odchylku přes hlavy. Vzorec je u~všech modelů týž: vrstva $0$ má nejvyšší labelovou a nejnižší kernelovou korelaci, střední vrstvy se chovají kernel-like (vysoká kernelová, nízká labelová korelace) a v~poslední vrstvě kernelová korelace opět klesá.}
    \label{fig:ch2_overview}
\end{figure}

Z~tabulky i obrázku vystupuje jasný a napříč modely stabilní vzorec.
\textbf{Vrstva 0} má u~všech vícevrstvých modelů nejvyšší labelovou korelaci ($|\text{corr}| = 0{,}22$ až $0{,}47$) a zároveň nejnižší kernelovou korelaci ($-0{,}10$ až $0{,}17$); u~osmivrstvého modelu je kernelová korelace první vrstvy dokonce mírně \textit{záporná}.
První vrstva tedy primárně \emph{čte labely} – shromažďuje informaci o~pozorovaných hodnotách, nikoli o~geometrii vstupů.
\textbf{Střední vrstvy} obraz obracejí: kernelová korelace vyskočí na $0{,}61$ až $0{,}89$ a labelová klesne na $0{,}01$ až $0{,}08$.
To je právě to \emph{kernel-like} chování, které bychom u~GP inference čekali – hlava porovnává vzdálenosti ve~vstupním prostoru.
\textbf{Poslední vrstva} pak kernelovou korelaci opět snižuje (na $0{,}27$ až $0{,}44$), zřejmě proto, že už neslouží k~porovnávání, ale k~sestavení konečné predikce.
\par

Zvláštní pozornost si zaslouží \textbf{jednovrstvý model}.
Jeho jediná vrstva má kernelovou korelaci $0{,}480$ – nejvyšší hodnotu v~celém sloupci vrstvy $0$ napříč modely.
To dává smysl: když má model k~dispozici jen jednu vrstvu, nemůže si dovolit ji „utratit“ na pouhé čtení labelů, musí v~ní stihnout jak label reading, tak kernelové porovnání současně.
Teprve s~přibývající hloubkou si model může dovolit specializaci, kde první vrstva labely jen posbírá a kernelovou práci deleguje na vrstvy pozdější.
Tuto specializaci na úrovni jednotlivých hlav ilustruje obrázek~\ref{fig:ch2_perhead}: hlavy uvnitř téže vrstvy nejsou zaměnitelné, každá pokrývá poněkud jiný aspekt vstupu.

\begin{figure}[htbp]
    \centering
    \includegraphics[width=\linewidth]{figures/ch2_label_kernel/fig_ch2_q1_8layer_corr_per_head.png}
    \caption{Osmivrstvý model, rozklad na jednotlivé hlavy. \textbf{Vlevo:} kernelová korelace, \textbf{vpravo:} labelová korelace, obojí pro každou z~$8$ hlav zvlášť napříč všemi $8$ vrstvami. Ve~vrstvě $0$ se hlavy rozcházejí – některé silně čtou labely (vpravo nahoře), jiné ne – zatímco ve~středních vrstvách se většina hlav sbíhá k~vysoké kernelové korelaci. Heterogenita hlav uvnitř vrstvy je motivací pro kauzální test v~oddílu~\ref{subsec:ch2-knockout}.}
    \label{fig:ch2_perhead}
\end{figure}

Stejnou heterogenitu lze nahlédnout i \emph{prostorově}.
Obrázek~\ref{fig:ch2_isoweight} zobrazuje, kam se dívá každá z~$8$ hlav jednovrstvého modelu: pro vybraný testovací bod $x^*$ vyneseme každý kontextový bod do roviny $(X_j, Y_j)$ a obarvíme ho jeho attention vahou.
Kdyby hlava prováděla kernelové porovnání, soustředila by váhu na body blízké $x^*$ (svislý pruh kolem $x^* = 0{,}59$), \emph{nezávisle} na jejich hodnotě $Y_j$.
Skutečnost je ale jiná: hlavy míří na \textbf{různé oblasti} roviny – Head~2 na body s~nízkým $X$ a vysokým $Y$ (vlevo nahoře), Head~1, 3, 5 a 7 naopak na body s~vysokým $X$ a nízkým $Y$ (vpravo dole) – a několik hlav (0, 4, 6) degeneruje na jediný bod (attention sink).
Protože kontextové body zde leží na klesající křivce GP realizace, znamená „dívat se do různých rohů“ dívat se na \textit{různé hodnoty} $Y_j$.
Attention tedy prokazatelně závisí i na labelu, nikoli jen na vzdálenosti ve~vstupním prostoru – což je prostorová obdoba nenulové labelové korelace z~tabulky~\ref{tab:ch2_corr_overview}.

\begin{figure}[htbp]
    \centering
    \includegraphics[width=\linewidth]{figures/ch2_label_kernel/fig_ch2_q1_iso_weight_grid.png}
    \caption{Iso-weight rozklad jednovrstvého modelu na $8$ hlav. Pro pevný testovací bod $x^* = 0{,}59$ (modrá přerušovaná čára) je každý kontextový bod vynesen do roviny $(X_j, Y_j)$ a obarven attention vahou dané hlavy (světlejší $=$ vyšší). Kernelové porovnání by koncentrovalo váhu kolem $x^*$ nezávisle na $Y_j$; místo toho různé hlavy míří do různých rohů roviny $(X,Y)$, tedy na různé hodnoty $Y_j$ – přímý vizuální doklad, že attention závisí i na labelu. Hlavy $0$, $4$, $6$ navíc degenerují na jediný bod (attention sink).}
    \label{fig:ch2_isoweight}
\end{figure}

Klíčový závěr korelační analýzy zní: \textit{nelze tvrdit, že PFN provádí čistý kernel smoothing}.
V~žádné vrstvě není labelová korelace nulová – model v~každém kroku míchá informaci o~vstupu $X$ i o~hodnotě $Y$.
Střední vrstvy sice kernel-like chování výrazně preferují, ale nikdy ho nedosahují v~čisté podobě.
Toto pozorování je plně v~souladu s~Naglerovým teoretickým výsledkem, ke~kterému se vracíme v~následujícím oddílu.

\clearpage
\subsection{Kauzální test důležitosti hlav: head-knockout}
\label{subsec:ch2-knockout}

Korelační analýza z~předchozího oddílu má zásadní omezení: \emph{pouze popisuje} attention vzorec.
Ignoruje value vektory, dopřednou síť (FFN) i reziduální tok, a proto z~ní nelze usuzovat na skutečnou funkční roli hlavy v~celém modelu.
Navíc Nagler dokázal (ve~své práci jako Věta~6.3), že softmax attention \textbf{vždy} míchá $X_j$ i $Y_j$ – „čistý kernel“ je tedy principiálně nedosažitelný a samotná (ne)korelace hlavy nám neřekne, zda je pro predikci důležitá.
Potřebujeme proto \textit{kauzální} metriku, která místo popisu měří skutečný dopad hlavy na výstup.
Slovem \textbf{kauzální důležitost} zde míníme něco jiného než korelaci: neptáme se, s~čím attention hlavy \emph{souvisí}, ale co se stane, když hlavu z~modelu \emph{skutečně odebereme}. Jde o~princip \emph{ablace} – zásah do modelu a změření jeho následku, analogický vyřazení komponenty v~neurovědě. Teprve dopad tohoto zásahu na predikci prozradí, zda byla komponenta pro výpočet nutná, nebo jen náhodně korelovala.
\par

Takovou metrikou je \textbf{head-knockout} (doslova „vyřazení hlavy“).
Připomeňme, že každá z~$8$ attention hlav vrstvy produkuje vlastní výstup, který se do reziduálního toku modelu přičítá přes naučenou výstupní projekci $W_{\text{out}}^{(h)}$.
Konkrétní hlavu proto můžeme \emph{chirurgicky odpojit} tak, že její výstupní projekci vynulujeme ($W_{\text{out}}^{(h)} = 0$) – ostatní hlavy, dopředné sítě i reziduální spoje přitom zůstanou nedotčené.
Poté modelem znovu proženeme tentýž vstup a změříme, o~kolik se zhorší shoda predikce s~GP posteriorem; rozdíl oproti neporušenému modelu je čistým příspěvkem právě té jedné hlavy.
Klíčové je použít \textbf{relativní} metriku:
$$\text{rel} = \frac{\text{MSE}(\text{knockout},\, \mu_{\text{GP}}) - \text{MSE}(\text{full},\, \mu_{\text{GP}})}{\text{MSE}(\text{full},\, \mu_{\text{GP}})}.$$
Absolutní $\Delta\text{MSE}$ by totiž byla \emph{zavádějící}: hluboké modely aproximují GP posterior s~triviálně malou baseline chybou (řádově $10^{-4}$), takže jejich absolutní zhoršení je malé z~principu, ať je hlava důležitá nebo ne. Relativní metrika (násobek vlastní baseline chyby modelu) měří skutečnou důležitost hlavy.
Velká hodnota $\text{rel}$ znamená, že hlava je pro aproximaci GP posterioru kauzálně nepostradatelná; hodnota blízká nule naopak značí, že model bez ní funguje prakticky stejně.
Test provádíme na $50$ instancích s~kontextem $n_{\text{ctx}} = 40$ a graf používá \textbf{per-model barevnou škálu} (jinak by hluboké modely s~malou baseline byly celé černé); výsledky shrnuje tabulka~\ref{tab:ch2_knockout} a obrázek~\ref{fig:ch2_knockout}.

\begin{table}[htbp]
\centering
\begin{tabular}{lccc}
\hline
\textbf{Model} & \textbf{Nejdůležitější (vrstva, hlava)} & \textbf{baseline MSE} & \textbf{max rel $\Delta$MSE} \\
\hline
$1$ vrstva & $(0,\, 7)$ & $0{,}01254$ & $\mathbf{147\times}$ \\
$2$ vrstvy & $(1,\, 5)$ & $0{,}00118$ & $21{,}4\times$ \\
$4$ vrstvy & $(1,\, 4)$ & $0{,}00015$ & $7{,}4\times$ \\
$6$ vrstev & $(2,\, 4)$ & $0{,}00195$ & $0{,}5\times$ \\
$8$ vrstev & $(1,\, 6)$ & $0{,}00008$ & $11{,}2\times$ \\
\hline
\end{tabular}
\caption{Kauzální důležitost hlav podle head-knockoutu, měřená \textbf{relativně} k~baseline chybě modelu. Uvedena je nejdůležitější hlava každého modelu, jeho baseline MSE vůči GP a maximální relativní zhoršení po vynulování té hlavy (např. $147\times$ znamená $147$násobek baseline chyby). Průměr přes $50$ instancí, $n_{\text{ctx}} = 40$.}
\label{tab:ch2_knockout}
\end{table}

\begin{figure}[htbp]
    \centering
    \includegraphics[width=\linewidth]{figures/ch2_label_kernel/fig_ch2_head_knockout_delta_mse.png}
    \caption{Head-knockout pro všech pět modelů, \textbf{per-model barevná škála}. Každá teplotní mapa zobrazuje \emph{relativní} zhoršení $\text{rel}=(\text{MSE}_{\text{knockout}}-\text{MSE}_{\text{full}})/\text{MSE}_{\text{full}}$ (světlejší $=$ důležitější) po vynulování hlavy na dané pozici (vrstva $\times$ hlava). Kauzálně důležité hlavy se koncentrují do \textbf{raných vrstev} (vrstva $0$–$1$) napříč všemi modely. Jednovrstvý model je nejkřehčí (až $147\times$ baseline), ale i hluboké modely nesou odebrání hlavy $7$–$21\times$ vlastní baseline chyby. Šestivrstvý model je výjimka (jen $0{,}5\times$, viz text).}
    \label{fig:ch2_knockout}
\end{figure}

Obrázek odhaluje jasnou strukturu: \textbf{kauzálně důležité hlavy se koncentrují do raných vrstev} – dominantní hlava je u~každého modelu ve~vrstvě $0$ nebo $1$.
To kauzálně potvrzuje hypotézu o~dělbě práce z~korelační analýzy: první vrstvy nesou hlavní výpočetní zátěž.
V~absolutních číslech hluboké modely knockout „ustojí“ (zůstávají blízko GP), avšak \textbf{relativně k~vlastní baseline chybě} zhorší odebrání jediné hlavy predikci pořád o~$7$ až $21$ násobek – hlavy tedy \emph{nejsou} bezvýznamné, jen individuálně nekritické.
Jednovrstvý model je nejkřehčí ($147\times$): s~jedinou vrstvou nese každá z~jejích osmi hlav nezastupitelný díl celého výpočtu.
\par

Výjimkou je \textbf{šestivrstvý model}, u~něhož je maximální relativní zhoršení jen $0{,}5\times$ a u~části hlav dokonce záporné – vynulování hlavy predikci mírně \emph{zlepší}.
Spolu s~jeho nezvykle vysokou baseline chybou ($0{,}00195$, vyšší než u~čtyř- i osmivrstvého modelu) to naznačuje, že tento konkrétní checkpoint je oproti ostatním hloubkám slabší, zřejmě méně dotrénovaný, a proto svou kapacitu plně nevyužívá.
\par

Podstatné je, že tento kauzální obraz korelační analýza \textit{ukázat nemohla}.
Hlava může s~kernelem korelovat silně, a přesto být nahraditelná; jiná může korelovat slabě, a přesto být kauzálně důležitá.
Head-knockout tak korelační obraz nejen doplňuje, ale v~jistém smyslu ho i koriguje – funkční roli hlavy určuje až její skutečný dopad na predikci, nikoli tvar jejího attention vzorce.

\clearpage
\subsection{Rozklad chyby a role hloubky}
\label{subsec:ch2-bias-variance}

Předchozí dva oddíly zkoumaly, \emph{co} vrstvy dělají.
Nyní se ptáme, \emph{kolik} jich je vlastně potřeba – tedy jak hloubka ovlivňuje samotnou přesnost aproximace GP posterioru.
K~tomu použijeme rozklad střední kvadratické chyby na dvě složky.
Pro každý model změříme MSE vůči GP posterioru při rostoucím počtu kontextových bodů $n \in \{5, 10, 20, 40, 80, 120\}$ ($200$ instancí na každé $n$) a naměřené hodnoty proložíme modelem
$$\text{MSE}(n) = \text{bias}^2 + \frac{c}{n}.$$
Zde $\text{bias}^2$ představuje \textit{ireducibilní} chybu – tu část, která nezmizí ani při nekonečném kontextu ($n \to \infty$) – zatímco člen $c/n$ je variance, která s~přibývajícími daty klesá.
Odhadnuté koeficienty jsou v~tabulce~\ref{tab:ch2_bias_variance}, proložení pak na obrázku~\ref{fig:ch2_msedecomp}.

\begin{table}[htbp]
\centering
\begin{tabular}{lccl}
\hline
\textbf{Model} & \textbf{$\text{bias}^2$} & \textbf{$c$ (variance)} & \textbf{Interpretace} \\
\hline
$1$ vrstva & $0{,}00484$ & $0{,}108$ & bias signifikantní \\
$2$ vrstvy & $-0{,}00121$ & $0{,}061$ & bias nesignifikantní \\
$4$ vrstvy & $-0{,}00098$ & $0{,}039$ & bias nesignifikantní \\
$6$ vrstev & $-0{,}00053$ & $0{,}025$ & bias nesignifikantní \\
$8$ vrstev & $-0{,}00065$ & $0{,}026$ & bias nesignifikantní \\
\hline
\end{tabular}
\caption{Rozklad MSE vůči GP posterioru na $\text{bias}^2 + c/n$. Jednovrstvý model má jako jediný signifikantně nenulový $\text{bias}^2$; od dvou vrstev výše je ireducibilní chyba statisticky neodlišitelná od nuly a hloubka už jen snižuje varianční člen $c$. Mírně \emph{záporné} odhady $\text{bias}^2$ u~hlubších modelů jsou artefaktem neváhovaného proložení – asymptota $\text{MSE}(n\to\infty)$ nemůže být fyzikálně záporná, takže tyto hodnoty prakticky znamenají nulový ireducibilní bias.}
\label{tab:ch2_bias_variance}
\end{table}

\begin{figure}[htbp]
    \centering
    \includegraphics[width=0.85\linewidth]{figures/ch2_label_kernel/fig_ch2_mse_vs_n_all_models.png}
    \caption{Rozklad MSE $=\text{bias}^2 + c/n$ pro všech pět modelů. Body jsou naměřené MSE (chybové úsečky $\pm 1$ std), přerušované křivky proložení. Jednovrstvý model (modrá) se s~rostoucím $n$ vyrovnává na~nenulové asymptotě $\text{bias}^2 = 0{,}0048$ (tečkovaná čára), zatímco hlubší modely klesají prakticky k~nule. Rozdíly mezi $4$, $6$ a $8$ vrstvami jsou už zanedbatelné.}
    \label{fig:ch2_msedecomp}
\end{figure}

Rozklad odhaluje dvě věci.
Za prvé, \textbf{jednovrstvý model má jako jediný signifikantně nenulový $\text{bias}^2 = 0{,}00484$}.
To je přímé empirické potvrzení Naglerova teoretického výsledku: jedna vrstva nestačí k~asymptotickému přiblížení GP posterioru, a proto v~ní zůstává ireducibilní chyba, kterou nelze odstranit žádným množstvím kontextových dat.
Od dvou vrstev výše je $\text{bias}^2 \approx 0$ – dvě vrstvy už stačí k~tomu, aby model při dostatku dat GP posterior dohnal libovolně přesně.
\par

Za druhé, hloubka nad rámec odstranění biasu působí především na \textbf{varianční člen $c$}, který monotónně klesá:
$$c_{1L} = 0{,}108 \;>\; c_{2L} = 0{,}061 \;>\; c_{4L} = 0{,}039 \;>\; c_{6L} = 0{,}025 \;\approx\; c_{8L} = 0{,}026.$$
Přidání vrstev tedy zrychluje konvergenci – hlubší model dosáhne téže přesnosti z~menšího kontextu.
Zlepšení se však sytí: mezi šesti a osmi vrstvami se $c$ prakticky nemění.
Z~praktického hlediska to znamená, že přibližně šest vrstev představuje bod, za kterým už další prohlubování sítě nepřináší téměř žádný zisk v~přesnosti.

\subsection{Shrnutí: co se která vrstva naučila}
\label{subsec:ch2-shrnuti}

Systematické měnění hloubky odhalilo, že PFN uvnitř sítě organizuje výpočet do zřetelné, byť nikoli dokonalé, dělby práce.
\par

Korelační analýza (pod Spearmanem) ukázala čistou \textbf{specializaci vrstev}: rané vrstvy čtou labely (nejvyšší $|\text{corr(label)}|$, nejnižší kernelová korelace), střední vrstvy se chovají kernel-like (kernelová korelace přes $0{,}8$) a poslední vrstva kernelové porovnávání opět opouští ve~prospěch sestavení predikce.
Zároveň však platí, že \textit{žádná} vrstva nedělá čistý kernel smoothing – v~každém kroku se míchá informace o~$X$ i o~$Y$, přesně jak předpovídá Naglerova věta o~softmax attention.
\par

Head-knockout tento popisný obraz doplnil \textbf{kauzálně}: kauzálně důležité hlavy se koncentrují do raných vrstev, jednovrstvý model je nejkřehčí (odebrání hlavy zhorší predikci až $147\times$ baseline), a i v~hlubokých modelech nese odebrání jediné hlavy $7$ až $21$ násobek vlastní baseline chyby – žádná hlava tedy není bezvýznamná, jen individuálně nekritická.
\par

Rozklad chyby vysvětlil, \textbf{proč modelu jedna vrstva nestačí}: jednovrstvý model má jako jediný nenulový $\text{bias}^2 = 0{,}0048$, zatímco od dvou vrstev výše je ireducibilní chyba nulová a hloubka už jen snižuje varianční člen (do nasycení kolem šesti vrstev).
Doplňkový test navíc ukázal, že identifikace korelační délky \textbf{neprobíhá} triviálním naladěním šířky attention v~první vrstvě (efektivní bandwidth vrstvy $0$ je vůči $\ell$ plochá), ale je rozprostřena hlouběji v~síti.
\par

Dohromady tato pozorování vykreslují PFN jako vrstvenou výpočetní architekturu, v~níž se – s~rostoucí hloubkou stále zřetelněji – dělí práce mezi sběr labelů, kernelové porovnávání a finální inferenci, aniž by kterákoli vrstva byla redukovatelná na prostý kernel smoothing.
